\section{Tevet--Berant Diversity-Eval: Human-Grounded Validation}\label{sec:tevet}

A diversity metric should be tested against human judgements on a standardized eval; \citet{tevet2020evaluating} provide both.
Their McDiv and ConTest datasets pair human-grounded diversity labels with a fixed comparison protocol (Spearman $\rho$ and OCA: \emph{optimal classification accuracy}, the best accuracy achievable by a one-dimensional threshold separating low- and high-diversity sets), the standard methodology for ranking diversity metrics against human judgements.
The response sets are written by Mechanical Turk workers, and for the McDiv\_nuggets subset the high/low label is fixed by the construction protocol itself: a worker writes five different continuations (the high-diversity set), then self-selects one of those continuations and paraphrases it five times (the low-diversity set; see Appendix~\ref{app:confound} for the full protocol and a construction-confound caveat).
We run $D_{Ca_n}$ through this evaluation pipeline.

\paragraph{Setup.} We use Qwen2.5-3B (base) with 50 permutations in completion format on the released no\_hds splits: \tevetMcDivSetCount{} McDiv sets, \tevetMcDivNugSetCount{} McDiv\_nuggets sets, and \tevetConTestSetCount{} ConTest sets, each with 5 worker-written responses per set and a binary high/low diversity label.
We follow Tevet and Berant in reporting Spearman $\rho$ and OCA, and additionally report ROC AUC (the natural metric for binary classification).

\begin{table*}[t]
\centering
\caption{ConTest results: Spearman $\rho$ and OCA between a set's diversity class and each metric score.
Baselines from Tevet's pre-computed metrics; $C \!\times\! a_n$ and $a_n$ are ours (Qwen2.5-3B, 50 permutations, per-byte).
Best result per task in bold.}
\label{tab:contest}
\small
% Generated by: scripts/analyze_c_ainf.py --run-tag qwen25_completion_v3
% Baseline source: diversity-eval/data/with_metrics/
\begin{tabular}{@{}l rr rr rr@{}}
\toprule
& \multicolumn{2}{c}{prompt\_gen} & \multicolumn{2}{c}{resp\_gen} & \multicolumn{2}{c}{story\_gen} \\
\cmidrule(lr){2-3} \cmidrule(lr){4-5} \cmidrule(lr){6-7}
Metric & $\rho$ & OCA & $\rho$ & OCA & $\rho$ & OCA \\
\midrule
\multicolumn{7}{@{}l}{\textit{ConTest (200, with\_hds)}} \\[2pt]
$C \!\times\! a_n$ (ours) & +0.584 & 0.785 & +0.391 & 0.668 & +0.686 & 0.828 \\
$a_n$ (ours) & +0.444 & 0.715 & +0.274 & 0.641 & +0.387 & 0.684 \\
SentBERT & \textbf{+0.682} & 0.815 & \textbf{+0.591} & \textbf{0.791} & \textbf{+0.770} & \textbf{0.896} \\
BERTsts & +0.646 & \textbf{0.820} & +0.463 & 0.714 & +0.601 & 0.780 \\
distinct-$n$ & +0.333 & 0.675 & +0.346 & 0.677 & +0.573 & 0.772 \\
\midrule
\multicolumn{7}{@{}l}{\textit{McDiv\_nuggets ($\sim$1K, no\_hds)}} \\[2pt]
$C \!\times\! a_n$ (ours) & +0.636 & 0.785 & +0.345 & 0.649 & +0.317 & 0.634 \\
$a_n$ (ours) & +0.487 & 0.705 & +0.225 & 0.619 & +0.124 & 0.557 \\
SentBERT & \textbf{+0.728} & \textbf{0.850} & \textbf{+0.532} & \textbf{0.758} & \textbf{+0.633} & \textbf{0.803} \\
BERTsts & +0.683 & 0.830 & +0.393 & 0.676 & +0.344 & 0.638 \\
distinct-$n$ & $-0.003$ & 0.514 & $-0.002$ & 0.507 & $-0.002$ & 0.510 \\
\midrule
\multicolumn{7}{@{}l}{\textit{McDiv (full, no\_hds, $\sim$2K)}} \\[2pt]
$C \!\times\! a_n$ (ours) & +0.729 & 0.846 & +0.500 & 0.724 & +0.523 & 0.717 \\
$a_n$ (ours) & +0.617 & 0.781 & +0.432 & 0.698 & +0.402 & 0.668 \\
SentBERT & \textbf{+0.796} & \textbf{0.897} & \textbf{+0.678} & \textbf{0.830} & \textbf{+0.753} & \textbf{0.867} \\
BERTsts & +0.780 & 0.893 & +0.614 & 0.781 & +0.571 & 0.740 \\
distinct-$n$ & +0.476 & 0.746 & +0.517 & 0.738 & +0.535 & 0.744 \\
\bottomrule
\end{tabular}

\end{table*}

\paragraph{Headline result.} On the binary tasks $D_{Ca_n}$ matches the strongest embedding baseline (Table~\ref{tab:contest}).
On McDiv prompt\_gen, $D_{Ca_n}$ reaches $\rho = \tevetMcDivPromptGenCxAnRho{}$ and OCA = \tevetMcDivPromptGenCxAnOCA{}, against SentBERT's $\rho = \tevetMcDivPromptGenSentBertRho{}$ and OCA = \tevetMcDivPromptGenSentBertOCA{}; the corresponding ROC AUC is \tevetMcDivPromptGenAUC{}.
Across the nine binary tasks, $D_{Ca_n}$'s OCA lags SentBERT by \tevetCxAnVsSentBertOCAMinPct\%--\tevetCxAnVsSentBertOCAMaxPct\%; it beats distinct-$n$ decisively on all three McDiv\_nuggets tasks (where distinct-$n$ is near chance) and trails it by 1--3 OCA points on three of the nine tasks where the labels are nearly shuffle-invariant.
The headline takeaway for this audience is that an information-theoretic, embedding-free metric reaches the same operating point as a sentence-embedding baseline on a human-grounded eval, with no embedding model, no reference corpus, and no human labels in the metric itself.

\paragraph{DecTest, included for completeness.} Tevet and Berant also release DecTest (Table~\ref{tab:dectest}), a diagnostic where the ``label'' is the sampling temperature used to generate each response set rather than a human judgement.
We include it for parity with the released benchmark only: temperature labels do not map onto creative-output diversity in any way that reflects the use cases in Section~\ref{sec:motivation}.
The raw $a_n$ achieves $\rho = \tevetDecTestAnPromptGenRho{}$ on prompt\_gen and $\rho = \tevetDecTestAnRespGenRho{}$ on resp\_gen, matching or exceeding all baselines.

\begin{table}[t]
\centering
\caption{DecTest results: Spearman $\rho$ between sampling temperature and each metric score (1000 samples, no\_hds).
Reported for parity with Tevet and Berant; the labels here are sampling temperatures rather than human judgements, so DecTest is not part of the framing for $D_{Ca_n}$.}
\label{tab:dectest}
\small
% Generated by: scripts/analyze_c_ainf.py --run-tag qwen25_completion_v3
% Baseline source: diversity-eval/data/with_metrics/
\begin{tabular}{@{}l rrr@{}}
\toprule
Metric & prompt\_gen & resp\_gen & story\_gen \\
\midrule
$C \!\times\! a_n$ (ours) & +0.847 & +0.771 & +0.763 \\
$a_n$ (ours) & \textbf{+0.932} & \textbf{+0.924} & \textbf{+0.779} \\
$C$ (ours) & $-0.727$ & $-0.813$ & $-0.547$ \\
distinct-$n$ & +0.917 & +0.894 & +0.758 \\
BERTScore & +0.878 & +0.874 & +0.694 \\
SentBERT & +0.747 & +0.801 & +0.645 \\
cos-sim & +0.873 & +0.895 & +0.712 \\
\bottomrule
\end{tabular}

\end{table}

\paragraph{Caveats.} The McDiv\_nuggets construction protocol introduces a confound: low-diversity sets are paraphrases of specific, dramatic endings that are intrinsically more surprising to the base model than typical high-diversity continuations.
Appendix~\ref{app:confound} documents the mechanism, the per-byte $a_1$ gap, and an evidence-by-length-bin breakdown showing the effect is content-driven not length-driven.
The $C$ weighting in $D_{Ca_n}$ helps on this benchmark partly for that confound-related reason; on settings without such a confound we expect $a_n$ to carry most of the signal.

